% Introduction (framing editor). Sources: research/00-brief.md (the request and the hypotheses H1--H10),
% CLAIMS.md §3 ("Answers to Hänni's requests"), and the section files. No theorem environments here.
\section{Introduction}\label{sec:intro}

In supervised learning a learner sees input--output pairs $(x_1,y_1),(x_2,y_2),\dots$ and predicts the output at a new input. For predicting sequences, Solomonoff induction (SI) is the standard ideal reference: a Bayes mixture over all computable sources whose cumulative log loss exceeds that of each computable source by at most a constant that depends on the source \citep{li2008introduction}. \citet{hanni2026notes} proposes the analogous object for supervised learning, \emph{Solomonoff function induction} (SFI): ``you think of the input-output pairs as being generated by picking a random program from this prior and generating using it, and so you update on seen pairs as you should''. This paper defines SFI carefully, proves its basic properties, checks the claims about it made in those notes, and works out examples.

\paragraph{Why SFI.} Supervised learning is central to machine learning, and SFI is a candidate for its ideal form. It is the Bayes predictor for the prior ``a random lower-semicomputable labeller''. On every sequence of pairs that some function could produce, its cumulative log loss exceeds that of every computable random labeller by at most that labeller's complexity plus a constant (\cref{thm:learn:dominance}). Since SFI never predicts the inputs, this holds for every input distribution, including adversarial and non-computable ones; with i.i.d.\ inputs it gives the Occam rate $O(K/n)$ for the averaged predictor, without a union bound (\cref{thm:learn:batch}). The analogy has limits, which we make precise: the guarantees are for log loss and for averaged predictors; a random function cannot produce fresh label noise on a repeated input; SFI is not computable; and SFI cannot use the order in which the data arrive (\cref{rem:learn:ideal}).

\paragraph{Relation to Solomonoff induction.} Up to constant factors, SI is SFI on the inputs $0,1,\dots,n-1$ (\cref{prop:setting:si}). In general, though, SFI sees the data as a set: each hypothesis answers each query without seeing the other inputs, and SFI's cumulative loss in the semimeasure convention does not depend on the order of the data (\cref{prop:setting:order}). SI applied to the interleaved sequence $x_1y_1x_2y_2\cdots$ sees the order and can do much better (\cref{sec:hist,sec:ex}). Giving SFI's hypotheses the list of previous inputs closes this gap: the result, \SFIhist{}, is the universal chronological semimeasure (\cref{cor:hist:H4}). With inputs as actions and labels as percepts this is the kind of object that underlies Hutter's AIXI \citep{hutter2000theory,hutter2005universal}; we have not checked the identification against Hutter's texts.

\paragraph{Relation to neural-network Bayes.} Bayesian learning with a neural network is function-space Bayes with the prior that the parameter prior induces on functions (\cref{prop:nn:functionspace}). \Citet{vaintrobhanni2026nnbayes} argue that under a strong weight prior this prior roughly tracks a circuit complexity of the function, so that NN Bayes behaves like a circuit version of SFI. \Cref{sec:comp,sec:nn} prove what can be proved in general about circuit SFI and about this link, and quote the slides' main claims as theirs.

\subsection{Contributions}\label{sec:intro:contrib}

\begin{itemize}[leftmargin=*]
\item \emph{Definitions and universality} (\cref{sec:setting}). Function machines read one random tape shared by all inputs. A universal function machine exists, and its law $\MF$ (SFI) dominates every machine law, including lower-semicomputable product laws, partial computable functions and laws of random functions with computable cylinder probabilities (\cref{thm:setting:universal,prop:setting:includes}). The one-point inequality alone defines a strictly larger class, whose universal element exceeds $\MF$ by a number of bits linear in the size of some data sets (\cref{thm:setting:weakvsmachine}). We prove invariance under the universal machine and under computable encodings (\cref{thm:setting:invariance}), coding-theorem bounds (\cref{thm:setting:coding}) and non-computability (\cref{prop:setting:computability}), and we separate the semimeasure convention for losses from the normalised one (\cref{prop:setting:normalise,rem:setting:conventions}).
\item \emph{Kolmogorov and semimeasure versions} (\cref{sec:versions}). The shortest consistent program $\KF$, the program mixture $W$, the product mixture $\Mprod$, SFI and the weak universal semimeasure agree up to $K(\dom D)+O(1)$ bits (\cref{thm:versions:chain}), but no two of $\KF$, $W$, $\MF$ agree up to a constant. There is no coding theorem for function programs: on some two-point data sets $\KF\ge\log d-2$ while $-\log W$ is $O(\log\log d)$ (\cref{thm:versions:nocoding,cor:versions:twopoint}). $W$ falls short of $\MF$ only by the padded complexity of the needed tape length (\cref{prop:versions:tape}). \Cref{tab:versions:compare} lists every comparison.
\item \emph{Learning theorems} (\cref{sec:learn}). Pathwise dominance (\cref{thm:learn:dominance}), merging (\cref{thm:learn:merging}), a mistake bound (\cref{cor:learn:mistakes}) and online-to-batch bounds (\cref{thm:learn:batch}). Under fresh label noise with repeated inputs, random-function SFI has infinite expected loss (\cref{prop:learn:freshnoise}); time stamps, frozen noise or a per-query mixture repair this (\cref{thm:learn:batch,prop:learn:perquery}). The expected risk of a Bayes predictor need not decrease with more data (\cref{prop:learn:curve}); per-step in-distribution error bounds are in \cref{cor:learn:perstep}.
\item \emph{History-aware prediction} (\cref{sec:hist}). On H\"anni's chain and skipped-string examples SFI's loss grows quadratically in the number of blocks when their lengths are $1,2,3,\dots$, in both conventions up to lower-order terms, while a constant-size set-based predictor has loss $0$, and every order-independent predictor pays the same lower bounds (\cref{thm:hist:chain,thm:hist:skipped,prop:hist:orderfree,prop:hist:normchain}). H\"anni's dominance theorem for computable reveal laws is true in the reading his text supports, under finiteness hypotheses; its proof sketch fails, and a stronger reading is false (\cref{prop:hist:sketchfails,cor:hist:thm1,thm:hist:randchain}). \SFIhist{} is the universal chronological semimeasure (\cref{thm:hist:representation,cor:hist:H4}); its exact relation to SI on the interleaved sequence is in \cref{thm:hist:si,prop:hist:inputloss}.
\item \emph{Worked examples} (\cref{sec:ex}). For a random cyclic permutation of $\{0,1\}^d$ revealed along its cycle, SFI pays $\log((N-1)!)\approx(d-\log e)2^d$ bits with high probability in the semimeasure convention, and in the normalised convention at least $\delta d-O(\log d)$ bits per step in expectation for the first $2^{(1-\delta)d}$ steps, while SI pays $O(1)$ in expectation and \SFIhist{} pays $O(1)$ on every instance, uniformly in $d$ (\cref{thm:ex:cycle,prop:ex:normalised,prop:ex:cycsi,prop:ex:cychist}). In Sam Eisenstat's coin example SFI pays $\frac{H(1/3)}2n\pm O(\log n)$ bits in the semimeasure convention, which is optimal among order-independent predictors up to $O(\log n)$, while SI pays at most $2.4$ bits plus a constant in expectation (\cref{thm:ex:expected,prop:ex:law,thm:ex:eisensi}). Positive examples are in \cref{sec:ex:positive}.
\item \emph{Out-of-distribution behaviour} (\cref{sec:ood}). An unseen simple property gives every simple behaviour on it a posterior probability bounded below, with no logarithmic term and independently of the amount of data (\cref{prop:ood:hanni}), and the bound depends on the universal machine (\cref{prop:ood:machine}). Over any sequence of distinct inputs, the total posterior error on a computable target $f$ is at most $(\ln2)(K(f)+c)$ (\cref{thm:ood:cumulative}). The naive converse fails, and a converse given the labelled sample holds (\cref{prop:ood:converse,prop:ood:labelled}).
\item \emph{Computable versions and networks} (\cref{sec:comp,sec:nn}). Time-bounded mixtures have constant regret against all predictors of a given running time (\cref{thm:comp:timebounded,prop:comp:round,cor:comp:poly}). For circuit Bayes, mistakes and circuit size are polynomially related on every ordering of the data (\cref{thm:comp:circuitbayes}), while simplest-consistent (Occam) learning on the same prior makes $2^{\Omega(\ell/\log\ell)}$ mistakes on some ordering for a target of description length $\ell$ (\cref{thm:comp:ladder}). For networks we prove the general facts about the induced prior and a separation between Gaussian-process limits and SFI on parities (\cref{prop:nn:upper,prop:nn:nolower,prop:nn:uniform,cor:nn:separation,rem:nn:oddparity}).
\item \emph{Steeper simplicity priors} (\cref{sec:steep}) summarises the companion paper \citep{claudeSFI2026steeper}: when a steeper penalty on program length keeps a learner from reproducing a teacher's systematic errors, and when it prefers a simple sampler to a model with filled-in random bits.
\end{itemize}

\paragraph{Corrections to earlier claims.} Several claims of \citet{hanni2026notes} are corrected here; the text says so where each appears. The proof sketch of the dominance theorem for reveal laws fails, although the theorem holds in its intended reading (\cref{prop:hist:sketchfails,cor:hist:thm1}). The two-sided bounds on out-of-distribution generalisation in the PDF draft are false as literally stated; the intended statement holds (\cref{prop:ood:draftfalse,thm:ood:twosided}). The converse asked about in the notes, ``no simple separating property implies good generalisation'', fails unless the bound may depend on $K(f)$ (\cref{prop:ood:converse,prop:ood:labelled}). The Kolmogorov and semimeasure versions, which the draft hoped to relate by a coding theorem, are not equivalent (\cref{rem:versions:answer}). The ``physicist's proof'' of few online mistakes is a proof for weighted majority, not for Occam learning (\cref{rem:comp:majority,thm:comp:ladder}). The e-mail argument for Eisenstat's example is correct in the semimeasure convention after two repairs (\cref{rem:ex:email}). \Cref{tab:intro:verdicts} lists the ten working hypotheses with which this study began and what became of each.

\begin{table}[!htbp]
\crefabbrev
\centering\footnotesize
\renewcommand{\arraystretch}{1.2}
\begin{tabularx}{\textwidth}{@{}>{\raggedright\arraybackslash}p{0.95cm}>{\raggedright\arraybackslash}p{4.2cm}>{\raggedright\arraybackslash}X@{}}
\toprule
& hypothesis (paraphrased) & verdict and where\\
\midrule
H1(a) & Universal lsc function semimeasures exist, equivalently given by function machines & Universal elements exist in both senses (\cref{thm:setting:universal,thm:setting:weakuniversal}). The equivalence is refuted for weak semimeasures, which can exceed $\MF$ by linearly many bits (\cref{thm:setting:weakvsmachine,rem:setting:h1a}); open for laws of random functions (\cref{q:setting:rpf}).\\
H1(b) & The program mixture loses about $K(\dom D)$, or a log factor via hashing, against SFI & All versions agree up to $K(\dom D)+O(1)$ (\cref{thm:versions:chain}). $W$ loses only a padded tape-length cost, with no log factor (\cref{prop:versions:tape}); the shortest program loses up to $\log d\pm O(\log\log d)$ on two points (\cref{thm:versions:nocoding,cor:versions:twopoint}); a log factor of the hashing kind describes $\KF$ (\cref{prop:versions:covering}).\\
H1(c) & The mixture of product laws is a third object & Strictly between $W$ and $\MF$, with unbounded gaps (\cref{prop:versions:Wgap,prop:versions:shared}); the per-query mixture is a fourth (\cref{def:versions:variants}).\\
H1 & SFI's cumulative loss does not depend on the order; input encoding matters little & Confirmed in the semimeasure convention (\cref{prop:setting:order}); normalised losses can depend on the order, open for $\MF$ (\cref{prop:setting:normalise}). Invariance up to $K(\alpha,\beta)+O(1)$ (\cref{thm:setting:invariance}).\\
H2 & Dominance over computable predictors that see only the input, on every input sequence & Confirmed for machine laws and consistent label sequences (\cref{thm:learn:dominance}); weak comparators need a $K(\dom D)$ term; open for general random-function laws (\cref{rem:learn:h2}).\\
H3 & H\"anni's counterexamples; his Theorem 1 for computable reveal laws & Counterexamples proved, also normalised (\cref{thm:hist:chain,thm:hist:skipped,prop:hist:normchain}). Theorem 1: proof sketch fails (\cref{prop:hist:sketchfails}); true in the intended reading with $K(n)+K(A)$ (\cref{cor:hist:thm1}); false in the oracle reading (\cref{thm:hist:randchain}).\\
H4 & With the list of previous inputs, SFI is the universal chronological semimeasure and as good as SI & First part confirmed (\cref{cor:hist:H4}); the set of previous inputs is not enough (\cref{prop:hist:strict}). Against SI's joint loss confirmed (\cref{thm:hist:si}); against its conditional label loss the difference is SI's loss on the inputs (\cref{prop:hist:inputloss}).\\
H5 & Cyclic permutation: high loss under SFI, low under SI & Confirmed (\cref{thm:ex:cycle,thm:ex:ksteps,prop:ex:cycsi}); normalised for $t\le2^{(1-\delta)d}$ (\cref{prop:ex:normalised}); over the full cycle normalised: open.\\
H6 & Eisenstat's example: SFI pays about $\frac{H(1/3)}2n$, SI a bounded amount & Confirmed up to $O(\log n)$ in the semimeasure convention (\cref{thm:ex:sandwich,thm:ex:expected,thm:ex:eisensi}); normalised: open (\cref{q:ex:normalised}).\\
H7 & For every input law and computable conditional law, expected regret $K(P)+O(1)$ and risk $K(P)/n$ & Confirmed for deterministic labels, time-stamped and frozen noise (\cref{thm:learn:batch}); refuted for fresh noise with repeated inputs (\cref{prop:learn:freshnoise}); holds for the per-query mixture (\cref{prop:learn:perquery}); last iterate open (\cref{rem:learn:lastiterate}).\\
H8 & An unseen simple property gives constant probability of bad behaviour; a converse? & Confirmed with no log term (\cref{prop:ood:hanni}); machine-dependent (\cref{prop:ood:machine}); naive converse refuted (\cref{prop:ood:converse}); converse given the labelled sample (\cref{prop:ood:labelled}).\\
H9 & A steeper simplicity prior with private randomness avoids a teacher's mistakes and splits model from randomness & Companion paper (\cref{sec:steep}): confirmed in a window of $\lambda/n$ for incompressible error rules; the extension to all error rules of that density and description length is refuted; $\lambda$ must grow linearly in $n$.\\
H10 & Time-bounded and circuit versions; NN Bayes approximates circuit SFI & Time-bounded regret (\cref{thm:comp:timebounded,prop:comp:round}); circuit mistakes versus size (\cref{thm:comp:circuitbayes}); Occam fails out of distribution (\cref{thm:comp:ladder}); NN link partly proved, partly claimed in the slides (\cref{sec:nn}).\\
\bottomrule
\end{tabularx}
\caption{The working hypotheses H1--H10, paraphrased, with verdicts. Most come from \citet{hanni2026notes,hanni2026request}; H1, H2 and H7 formalise natural expectations. Lower bounds are in the semimeasure convention unless stated.}
\label{tab:intro:verdicts}
\end{table}

\paragraph{Credit.} The notion of SFI, many of the claims examined here and most of the examples come from \citet{hanni2026notes} and \citet{hanni2026request}: the dominance template for samplers that see only the input, the remark on input encodings, the chain and skipped-string examples, the dominance theorem for reveal laws, the conjecture on hypotheses that see the previous inputs, the cyclic-permutation example, the out-of-distribution proposition and examples, and the time-bounded and circuit versions. The coin example of \cref{sec:ex:eisenstat} is due to Sam Eisenstat, as reported in \citet{hanni2026request}. The material on neural-network Bayes, including the conjunction ladder of \cref{thm:comp:ladder}, follows the slides of \citet{vaintrobhanni2026nnbayes}. Hutter's universal chronological semimeasure and his treatment of supervised learning as a problem class for AIXI \citep{hutter2000theory,hutter2005universal,hutter2007aixi} are prior work in substance for \cref{sec:hist}; we have checked them only at the level of abstracts and tables of contents (\cref{rem:learn:hutterEX,sec:related}).

\paragraph{What is not achieved.} Several natural questions stay open: whether $\MF$ dominates every lower-semicomputable law of a random partial function (\cref{q:setting:rpf}); the exact size of the gap between $W$ and $\MF$ (\cref{q:versions:open}); several lower bounds in the normalised convention (\cref{q:setting:normorder,q:ex:normalised}); whether the $K(n)+K(A)$ term in the dominance theorem for reveal laws can be removed (\cref{rem:hist:thm1open}); an $O(K/n)$ bound for the last iterate rather than the average (\cref{rem:learn:lastiterate,conj:learn:loo}); efficient time-bounded learning against samplers (\cref{q:comp:samplers}); and a two-sided quantitative bridge between NN Bayes and circuit SFI, of which one half is proved for uniform circuit priors (\cref{prop:nn:uniform,q:nn:bridges}). The main theorems of the slides are not checked here. The constants in our bounds depend on the universal machine and can be large, and SFI itself is not computable, so the results describe an ideal, not an algorithm. Much of the literature was seen only through abstracts and secondary sources (\cref{sec:related}).

\paragraph{Reader's guide.} \Cref{sec:setting} is the base for everything else. \Cref{sec:versions} answers the question of the Kolmogorov and semimeasure versions and can be skipped on a first reading. \Cref{sec:learn} gives the learning theorems. \Cref{sec:hist,sec:ex} show where SFI and SI differ and how history-aware hypotheses remove the difference. \Cref{sec:ood} treats out-of-distribution behaviour, \cref{sec:comp,sec:nn} computable versions and networks, \cref{sec:steep} the companion paper, and \cref{sec:related,sec:discussion} related work and open problems. Each main section states its results with short proof ideas; full proofs are in \cref{app:setting,app:versions,app:learn,app:hist,app:ex,app:ood,app:comp,app:nn}. Two conventions for losses recur: the \emph{semimeasure convention} scores SFI's raw predictive sub-probabilities, and the \emph{normalised convention} rescales them to sum to one at each step. Upper bounds in the first hold in both; every lower bound says which convention it is for (\cref{rem:setting:conventions}). Logarithms are to base~$2$, losses are in bits, and $c$ denotes a constant that depends only on the fixed universal machines and encodings.
